% IV. Програмна реализация.
% Попълва се с код и коментар след реализацията. Тук стои само планираната структура.
\section{Програмна реализация}
\label{sec:implementation}

\subsection{Организация на изходния текст}

Изходният текст е разделен на библиотека с протоколното ядро, библиотека със слоя за
вход, изход и време, и две изпълними програми, които ги свързват: демонът и
симулаторът. Разделението е принудително, а не по вкус: ако ядрото е отделна библиотека
без зависимост към сокети, нарушение на границата се проявява като грешка при свързване,
а не като незабелязана зависимост. Изграждането е с CMake, а тестовете използват
GoogleTest.

\subsection{Автомат за състоянията на съседите}

Съседството преминава през последователност от състояния, в която двата възела
установяват двупосочна видимост и след това синхронизират базите си, преди да обявят
връзката за използваема~\cite{rfc2328}. Реализацията описва автомата като явна таблица
от преходи, вместо като вложени условни оператори, за да може всеки преход да бъде
покрит с отделен тест.
\TODO{да се приложи таблицата на преходите след реализацията и да се посочи покритието}

\subsection{Разпространение на обявленията}

Обявлението за състояние на връзка носи подател, пореден номер, възраст и списък от
съседи с цени на връзките. Приемането сравнява поредния номер с този на записаното
обявление от същия подател, записва по-новото и препраща. Проверката за по-нов номер е
критична: без нея лавинното разпространение не спира.
\TODO{да се опише поведението при обръщане на поредния номер и при остаряло обявление}

\subsection{Изчисляване и инсталиране на маршрути}

Изчисляването е реализация на алгоритъма на Дийкстра~\cite{dijkstra1959} върху графа,
реконструиран от базата, с опашка с приоритет. Резултатът е следваща стъпка за всяка
достижима цел. Инсталирането е отделна стъпка, за да остане измерим интервалът между
изчисляването и влизането на маршрута в сила.
\TODO{да се опише как маршрутът се инсталира в маршрутната таблица на операционната
система в режим на демон и какви права изисква това на всяка от трите платформи}

\subsection{Съвместимост с еталонната реализация}

Съвместимостта на формата на пакетите се проверява чрез обмен с еталонна реализация в
изолирана среда. Проверява се, че еталонът установява съседство с демона и че
изчислените от двете страни маршрути съвпадат за една и съща топология.
\TODO{да се опише средата за проверката и да се приложи резултатът}
